%--------------------------------------------------------------------%
% Page          : Abstract (Indonesian)
%--------------------------------------------------------------------%

\clearpage

\addcontentsline{toc}{chapter}{ABSTRAK}

\section*{\large \centering \MakeUppercase 
    {Abstrak}
}

\begin{center}

    \justifying \normalsize

    \textbf{Pengambilan Keputusan dan Pengaturan Skalabilitas Elastis pada Lingkungan Cloud yang Heterogen dengan Berbasiskan pada Prediksi Workload}

    \vspace{0.5cm}

    \textbf{Muhammad Zahid Masruri / NRP 6025222041}

    \vspace{0.5cm}

    Arsitektur \emph{cloud} modern menghadapi tantangan \emph{trade-off} antara skalabilitas dan efisiensi biaya. \emph{Kubernetes} menyediakan kontrol infrastruktur yang baik namun lambat merespons lonjakan lalu lintas, sementara komputasi \emph{serverless} menawarkan elastisitas tinggi dengan biaya per-\emph{request} yang lebih mahal. Penelitian ini merancang, mengimplementasikan, dan mengevaluasi arsitektur \emph{hybrid} yang mengintegrasikan kluster \emph{Kubernetes} dan \emph{serverless} dengan prediksi beban kerja berbasis \emph{Gated Recurrent Unit} (GRU) untuk distribusi lalu lintas yang proaktif.

    Sistem yang diusulkan mengimplementasikan pengontrol routing V3 berbasis kapasitas (\emph{capacity-driven}) yang menggeser lalu lintas secara dinamis antara \emph{backend Kubernetes} dan \emph{serverless} berdasarkan pemantauan \emph{tail latency}. Mekanisme \emph{proactive routing} menggunakan ekstrapolasi tren beban kerja aktual (5 langkah $\times$ 15 detik = 75 detik ke depan) yang dipicu oleh tingkat kepercayaan GRU. Evaluasi dilakukan melalui 30 eksperimen terkontrol yang direplikasi pada empat skenario: K8s+HPA (\emph{baseline}), \emph{Serverless}-only, \emph{Hybrid-reactive}, dan \emph{Hybrid-predictive}, dengan beban \emph{trace} ClarkNet (40 tahap $\times$ 30 detik = 20 menit, RPS 22--164).

    Hasil menunjukkan S4 (\emph{Hybrid-predictive}) mencapai latensi p99 = 188\,ms dengan 702 pelanggaran SLO per \emph{run}, 75\% lebih sedikit dibandingkan S3 (\emph{Hybrid-reactive}) dengan 2.805 pelanggaran (\emph{Cohen's d} = 1,21, \emph{large effect}). S4 juga 20\% lebih murah daripada S1 (\$149/bulan vs \$187/bulan) dengan tingkat kepatuhan SLO 98,4\%. Model GRU mencapai 400/400 prediksi sukses dengan tingkat kepercayaan 0,72--0,82 dan latensi inferensi 10--13\,ms. Mekanisme \emph{proactive routing} memicu 9 aksi PREDICTIVE per \emph{run}, memulai distribusi ke \emph{serverless} pada 65\% kapasitas K8s.

    \vspace{0.5cm}

    \textbf{Kata Kunci:} \emph{Cloud Computing, Kubernetes, Serverless, Prediksi Beban Kerja, GRU, Distribusi Lalu Lintas, Arsitektur Hybrid}

\end{center}

\clearpage
